\documentclass[12pt]{hw}


%====================================================================
% IMPORTANT INFORMATION TO FILL OUT:
%====================================================================
\usepackage{amsthm}
\newtheorem{definition}{Definition}
\newcommand{\getr}{\ensuremath{{\overset{\$}{\leftarrow}}}\xspace}
\newcommand{\Gen}{\ensuremath{{\sf Gen}}\xspace}
\newcommand{\D}{\mathcal{D}}
% ENTER YOUR NAME HERE:
\renewcommand{\student}{ YOUR NAME HERE (uniqname here) }

% LIST ALL COLLABORATORS HERE:
\renewcommand{\collaborators}{Name/Exercise Number(s), Name/Exercise Numbers (s)}

%====================================================================
%====================================================================

\setcounter{sheet}{3} % Homework number here

\begin{document}

\maketitle

\begin{center} Due: Oct 10, 2024 (at 11:59 PM US Eastern Time) \\[10pt]
  \hrule height 1pt
\end{center}

\toggletrue{answers}
% \togglefalse{answers}

\setcounter{MaxMatrixCols}{20} \setlength\parindent{0pt}

% =============================================================================
% =============================================================================


% ===============================================
% ===============================================


% =============================================================================


\begin{exercise}

  \textbf{(20 points)} %
  Let $F\colon \bit^n \times \bit^n \rightarrow \bit^n$ be pseudorandom function.
  Show that each of the following MACs is insecure, even if used to
  authenticate fixed-length
  messages. (In each case $\skcgen$ outputs a uniform $k \in \bit^n$,
  $n$ being even, and $\langle i \rangle$ denotes an $n/2$-bit binary
  encoding of the integer $i$).
  \begin{enumerate}[(a)]
    \item To authenticate a message $m = m_1\| \cdots \| m_\ell$, where
          $m_i \in \bit^n$, compute \\
          $t := F_k(m_1)\oplus \cdots \oplus F_k(m_\ell)$.
    \item To authenticate a message $m = m_1 \| \cdots \| m_\ell$, where
          $m_i \in \bit^{n/2}$, compute \\
          $t := F_k(\langle 1 \rangle \| m_1) \oplus \cdots \oplus
            F_k(\langle \ell \rangle \| m_\ell)$.
  \end{enumerate}
\end{exercise}

\begin{answer}
  \begin{enumerate}[(a)]
    \item First part.

    \item Second part.

  \end{enumerate}
\end{answer}


% =============================================================================


\begin{exercise}
  \textbf{(15 points)}
  Let $F\colon \bit^n \times \bit^n \rightarrow \bit^n$ be a PRF. 
  Consider the following\\ $\mac=(\macgen,\mactag,\macver)$ given by
  \begin{itemize}
    \item $\macgen(1^n)$: outputs a uniformly random key $k\gets\bit^n$.
    \item $\mactag(k,m)$: given a message $m=m_1\|m_2\in\bit^{2n}$, 
      where $|m_1|=|m_2|=n$, output the tag $t:=F_k(m_1)\| F_k(F_k(m_2))$.
    \item $\macver(k,m,t)$: check whether $\mactag(k,m)=t$.
  \end{itemize}
  Is this MAC secure (existential unforgeabilty under chosen message attack)? Prove your answer either by giving a formal reduction
  to $F$ being a PRF or by constructing a forger with non-negligible advantage
  in the security game against $\mac$. 
\end{exercise}

\begin{answer} 
Write answer here.
\end{answer}


% =============================================================================
\begin{exercise} \textbf{(20 points)}
  For each of the following hash functions $H$, determine whether it
  is \emph{necessarily} collision resistant. If it is, describe and
  analyze an appropriate reduction. If it's not, describe and analyze
  a feasible attacker against the collision resistance of $H$.
  
  \begin{enumerate}[(a)]
  \item $H^{s}(x) = H_1^{s}(x) \oplus H_2^{s}(x)$, where $\{H^s_1\}, \{H^s_2\}$ are both
    collision-resistant hash functions families.
  
  \item $H^s(x) = \widetilde{H}^s(\widetilde{H}^s(x))$, where $\{\widetilde{H}^s\}$ is a collision-resistant hash
    function family (that can support arbitrary length inputs).
  \end{enumerate}

\end{exercise}
\begin{answer}
Write Solution Here.
\end{answer}
\begin{exercise}
\textbf{(15 points)}
Recall the following definitions for MAC schemes.

\begin{definition}[\bf{Existentially Unforgeable under Chosen Message}]
    A MAC scheme $(\mathsf{Gen}, \mathsf{Mac}, \mathsf{Vrfy})$ is \emph{existentially unforgable under a chosen-message attack} if and only for every PPT adversary $\mathcal{A}$, 
    there is a negligible function $\mathsf{negl}(\cdot)$ such that 
for all $n \in \mathbb{N}$, $\Pr[\mathsf{Expt}_{\mathsf{EU}}^{\mathcal{A}}(1^n) = 1] \leq \mathsf{negl}(n)$.
    \begin{mdframed}
    
    ${\bf \mathsf{Expt}_{\mathsf{EU}}^{\mathcal{A}}(1^n)}$.
    
    \begin{itemize}
        \item $k \getr \Gen(1^n)$; 
        \item $(m, t) \leftarrow \mathcal{A}^{{\sf MAC}_k(\cdot)}(1^n)$; (The adversary has oracle access to the MAC function parameterized with k to get tags of messages)
        \item Output 1 if ${\sf Vrfy}(k,m,t) = 1$ and $m$ not queried by $\mathcal{A}$.
    \end{itemize}
    \end{mdframed}
\end{definition}
In words, the adversary is able to make queries to get tags for a polynomial amount of messages and then make a message tag pair itself that is not of a message previously queried.
\begin{definition}[\bf{Strongly Existentially Unforgeable under Chosen Message}]
    A MAC scheme $(\mathbf{Gen}, \mathbf{Mac}, \mathbf{Vrfy})$ is \emph{strongly existentially unforgable under a chosen-message attack} if and only for every PPT adversary $\mathcal{A}$, 
    there is a negligible function $\mathsf{negl}(\cdot)$ such that 
for all $n\in \mathbb{N}$, $\Pr[\mathsf{Expt}_{\mathsf{SEU}}^{\mathcal{A}}(1^n) = 1] \leq \mathsf{negl}(n)$.
    
    \begin{mdframed}
    
    ${\bf \mathsf{Expt}_{\mathsf{SEU}}^{\mathcal{A}}(1^n)}$.
    
    \begin{itemize}
        \item $k \getr \Gen(1^n)$;
        \item $(m, t) \leftarrow \mathcal{A}^{{\sf MAC}_k(\cdot)}(1^n)$; (The adversary has oracle access to the MAC function parameterized with k to get tags of messages)
        \item Output 1 if ${\sf Vrfy}(k,m,t) = 1$ and $(m,t)$ is not a pair seen in the set of queries made by $\mathcal{A}$.
    \end{itemize}
    \end{mdframed}
\end{definition}
In words, the adversary is able to make queries to get tags for a polynomial amount of messages and then make a message tag pair itself that \textbf{can be a message previously queried, but the tag cannot be a tag associated with that message from a previous query.}\\

Prove that if a MAC scheme $(\mathsf{Gen}, \mathsf{Mac}, \mathsf{Vrfy})$ is strongly existentially unforgable under a chosen-message attack, then it is existentially unforgable under a chosen-message attack.
\end{exercise}
\begin{answer}
Write Solution Here.
\end{answer}
\begin{exercise} {\bf (30 points)}
    Let $H\colon \bit^\ell \to \bit^n$ be a collision-resistant hash function,
    where $\ell > n$, and $F\colon \keyspace \times \bit^n \to \bit^n$ be a PRF.
    Show that the function $F'\colon \keyspace \times \bit^\ell \to \bit^n$
    given by $F'_k(x) := F_k(H(x))$ is a PRF.  In practice, this construction
    can produce a PRF with arbitrary-length inputs.
    \smallskip

    \emph{Hint: Consider the similarities and differences between this
    construction and the hash-and-tag construction that we analyzed in lecture.}
\end{exercise}
    
\begin{answer}
Write Solution Here.
\end{answer}
% =============================================================================
\end{document}
